\documentclass[
  doc,
  floatsintext,
  longtable,
  a4paper,
  nolmodern,
  notxfonts,
  notimes,
  12pt,
  colorlinks=true,linkcolor=blue,citecolor=blue,urlcolor=blue]{apa7}

\usepackage{amsmath}
\usepackage{amssymb}

\geometry{inner=1in, outer=1in}
\fancyhfoffset[LE,RO]{0cm}


\usepackage[bidi=default]{babel}
\babelprovide[main,import]{spanish}


% get rid of language-specific shorthands (see #6817):
\let\LanguageShortHands\languageshorthands
\def\languageshorthands#1{}

\RequirePackage{longtable}
\RequirePackage{threeparttablex}

\makeatletter
\renewcommand{\paragraph}{\@startsection{paragraph}{4}{\parindent}%
	{0\baselineskip \@plus 0.2ex \@minus 0.2ex}%
	{-.5em}%
	{\normalfont\normalsize\bfseries\typesectitle}}

\renewcommand{\subparagraph}[1]{\@startsection{subparagraph}{5}{0.5em}%
	{0\baselineskip \@plus 0.2ex \@minus 0.2ex}%
	{-\z@\relax}%
	{\normalfont\normalsize\bfseries\itshape\hspace{\parindent}{#1}\textit{\addperi}}{\relax}}
\makeatother




\usepackage{longtable, booktabs, multirow, multicol, colortbl, hhline, caption, array, float, xpatch}
\usepackage{subcaption}


\renewcommand\thesubfigure{\Alph{subfigure}}
\setcounter{topnumber}{2}
\setcounter{bottomnumber}{2}
\setcounter{totalnumber}{4}
\renewcommand{\topfraction}{0.85}
\renewcommand{\bottomfraction}{0.85}
\renewcommand{\textfraction}{0.15}
\renewcommand{\floatpagefraction}{0.7}

\usepackage{tcolorbox}
\tcbuselibrary{listings,theorems, breakable, skins}
\usepackage{fontawesome5}

\definecolor{quarto-callout-color}{HTML}{909090}
\definecolor{quarto-callout-note-color}{HTML}{0758E5}
\definecolor{quarto-callout-important-color}{HTML}{CC1914}
\definecolor{quarto-callout-warning-color}{HTML}{EB9113}
\definecolor{quarto-callout-tip-color}{HTML}{00A047}
\definecolor{quarto-callout-caution-color}{HTML}{FC5300}
\definecolor{quarto-callout-color-frame}{HTML}{ACACAC}
\definecolor{quarto-callout-note-color-frame}{HTML}{4582EC}
\definecolor{quarto-callout-important-color-frame}{HTML}{D9534F}
\definecolor{quarto-callout-warning-color-frame}{HTML}{F0AD4E}
\definecolor{quarto-callout-tip-color-frame}{HTML}{02B875}
\definecolor{quarto-callout-caution-color-frame}{HTML}{FD7E14}

%\newlength\Oldarrayrulewidth
%\newlength\Oldtabcolsep


\usepackage{hyperref}



\usepackage{color}
\usepackage{fancyvrb}
\newcommand{\VerbBar}{|}
\newcommand{\VERB}{\Verb[commandchars=\\\{\}]}
\DefineVerbatimEnvironment{Highlighting}{Verbatim}{commandchars=\\\{\}}
% Add ',fontsize=\small' for more characters per line
\usepackage{framed}
\definecolor{shadecolor}{RGB}{241,243,245}
\newenvironment{Shaded}{\begin{snugshade}}{\end{snugshade}}
\newcommand{\AlertTok}[1]{\textcolor[rgb]{0.68,0.00,0.00}{#1}}
\newcommand{\AnnotationTok}[1]{\textcolor[rgb]{0.37,0.37,0.37}{#1}}
\newcommand{\AttributeTok}[1]{\textcolor[rgb]{0.40,0.45,0.13}{#1}}
\newcommand{\BaseNTok}[1]{\textcolor[rgb]{0.68,0.00,0.00}{#1}}
\newcommand{\BuiltInTok}[1]{\textcolor[rgb]{0.00,0.23,0.31}{#1}}
\newcommand{\CharTok}[1]{\textcolor[rgb]{0.13,0.47,0.30}{#1}}
\newcommand{\CommentTok}[1]{\textcolor[rgb]{0.37,0.37,0.37}{#1}}
\newcommand{\CommentVarTok}[1]{\textcolor[rgb]{0.37,0.37,0.37}{\textit{#1}}}
\newcommand{\ConstantTok}[1]{\textcolor[rgb]{0.56,0.35,0.01}{#1}}
\newcommand{\ControlFlowTok}[1]{\textcolor[rgb]{0.00,0.23,0.31}{\textbf{#1}}}
\newcommand{\DataTypeTok}[1]{\textcolor[rgb]{0.68,0.00,0.00}{#1}}
\newcommand{\DecValTok}[1]{\textcolor[rgb]{0.68,0.00,0.00}{#1}}
\newcommand{\DocumentationTok}[1]{\textcolor[rgb]{0.37,0.37,0.37}{\textit{#1}}}
\newcommand{\ErrorTok}[1]{\textcolor[rgb]{0.68,0.00,0.00}{#1}}
\newcommand{\ExtensionTok}[1]{\textcolor[rgb]{0.00,0.23,0.31}{#1}}
\newcommand{\FloatTok}[1]{\textcolor[rgb]{0.68,0.00,0.00}{#1}}
\newcommand{\FunctionTok}[1]{\textcolor[rgb]{0.28,0.35,0.67}{#1}}
\newcommand{\ImportTok}[1]{\textcolor[rgb]{0.00,0.46,0.62}{#1}}
\newcommand{\InformationTok}[1]{\textcolor[rgb]{0.37,0.37,0.37}{#1}}
\newcommand{\KeywordTok}[1]{\textcolor[rgb]{0.00,0.23,0.31}{\textbf{#1}}}
\newcommand{\NormalTok}[1]{\textcolor[rgb]{0.00,0.23,0.31}{#1}}
\newcommand{\OperatorTok}[1]{\textcolor[rgb]{0.37,0.37,0.37}{#1}}
\newcommand{\OtherTok}[1]{\textcolor[rgb]{0.00,0.23,0.31}{#1}}
\newcommand{\PreprocessorTok}[1]{\textcolor[rgb]{0.68,0.00,0.00}{#1}}
\newcommand{\RegionMarkerTok}[1]{\textcolor[rgb]{0.00,0.23,0.31}{#1}}
\newcommand{\SpecialCharTok}[1]{\textcolor[rgb]{0.37,0.37,0.37}{#1}}
\newcommand{\SpecialStringTok}[1]{\textcolor[rgb]{0.13,0.47,0.30}{#1}}
\newcommand{\StringTok}[1]{\textcolor[rgb]{0.13,0.47,0.30}{#1}}
\newcommand{\VariableTok}[1]{\textcolor[rgb]{0.07,0.07,0.07}{#1}}
\newcommand{\VerbatimStringTok}[1]{\textcolor[rgb]{0.13,0.47,0.30}{#1}}
\newcommand{\WarningTok}[1]{\textcolor[rgb]{0.37,0.37,0.37}{\textit{#1}}}

\providecommand{\tightlist}{%
  \setlength{\itemsep}{0pt}\setlength{\parskip}{0pt}}
\usepackage{longtable,booktabs,array}
\newcounter{none} % for unnumbered tables
\usepackage{calc} % for calculating minipage widths
% Correct order of tables after \paragraph or \subparagraph
\usepackage{etoolbox}
\makeatletter
\patchcmd\longtable{\par}{\if@noskipsec\mbox{}\fi\par}{}{}
\makeatother
% Allow footnotes in longtable head/foot
\IfFileExists{footnotehyper.sty}{\usepackage{footnotehyper}}{\usepackage{footnote}}
\makesavenoteenv{longtable}

\usepackage{graphicx}
\makeatletter
\newsavebox\pandoc@box
\newcommand*\pandocbounded[1]{% scales image to fit in text height/width
  \sbox\pandoc@box{#1}%
  \Gscale@div\@tempa{\textheight}{\dimexpr\ht\pandoc@box+\dp\pandoc@box\relax}%
  \Gscale@div\@tempb{\linewidth}{\wd\pandoc@box}%
  \ifdim\@tempb\p@<\@tempa\p@\let\@tempa\@tempb\fi% select the smaller of both
  \ifdim\@tempa\p@<\p@\scalebox{\@tempa}{\usebox\pandoc@box}%
  \else\usebox{\pandoc@box}%
  \fi%
}
% Set default figure placement to htbp
\def\fps@figure{htbp}
\makeatother







\usepackage{newtx}

\defaultfontfeatures{Scale=MatchLowercase}
\defaultfontfeatures[\rmfamily]{Ligatures=TeX,Scale=1}





\title{Tipos de datos}


\shorttitle{TIPOS DE DATOS}


\usepackage{etoolbox}



\ccoppy{\textcopyright~2026}



\author{Edison Achalma}



\affiliation{
{Escuela Profesional de Economía, Universidad Nacional de San Cristóbal
de Huamanga}}




\leftheader{Achalma}

\date{2026-10-20}


\abstract{Este abstract será actualizado una vez que se complete el
contenido final del artículo. }



\authornote{\par{\addORCIDlink{Edison Achalma}{0000-0001-6996-3364}} 
\par{ }
\par{   El autor no tiene conflictos de interés que revelar.    Los
roles de autor se clasificaron utilizando la taxonomía de roles de
colaborador (CRediT; https://credit.niso.org/) de la siguiente
manera:  Edison Achalma:   conceptualización, metodología, análisis
formal, investigación, recursos, curación de
datos, redacción, visualización, supervisión, administración del
proyecto}
\par{La correspondencia relativa a este artículo debe dirigirse a Edison
Achalma, Escuela Profesional de Economía, Universidad Nacional de San
Cristóbal de Huamanga, Portal Independencia N°
57, Ayacucho, AYA 5001, Perú, Email: \href{mailto:elmer.achalma.09@unsch.edu.pe}{elmer.achalma.09@unsch.edu.pe}}
}

\makeatletter
\let\endoldlt\endlongtable
\def\endlongtable{
\hline
\endoldlt
}
\makeatother

\urlstyle{same}



\hypersetup{
  pdftitle={Tipos de datos},
  pdfauthor={Autor}
}
\makeatletter
\@ifpackageloaded{caption}{}{\usepackage{caption}}
\AtBeginDocument{%
\ifdefined\contentsname
  \renewcommand*\contentsname{Tabla de contenidos}
\else
  \newcommand\contentsname{Tabla de contenidos}
\fi
\ifdefined\listfigurename
  \renewcommand*\listfigurename{Lista de Figuras}
\else
  \newcommand\listfigurename{Lista de Figuras}
\fi
\ifdefined\listtablename
  \renewcommand*\listtablename{Lista de Tablas}
\else
  \newcommand\listtablename{Lista de Tablas}
\fi
\ifdefined\figurename
  \renewcommand*\figurename{Figura}
\else
  \newcommand\figurename{Figura}
\fi
\ifdefined\tablename
  \renewcommand*\tablename{Tabla}
\else
  \newcommand\tablename{Tabla}
\fi
}
\@ifpackageloaded{float}{}{\usepackage{float}}
\floatstyle{ruled}
\@ifundefined{c@chapter}{\newfloat{codelisting}{h}{lop}}{\newfloat{codelisting}{h}{lop}[chapter]}
\floatname{codelisting}{Listado}
\newcommand*\listoflistings{\listof{codelisting}{Lista de Listados}}
\makeatother
\makeatletter
\makeatother
\makeatletter
\@ifpackageloaded{caption}{}{\usepackage{caption}}
\@ifpackageloaded{subcaption}{}{\usepackage{subcaption}}
\makeatother
\makeatletter
\@ifpackageloaded{fontawesome5}{}{\usepackage{fontawesome5}}
\makeatother

% From https://tex.stackexchange.com/a/645996/211326
%%% apa7 doesn't want to add appendix section titles in the toc
%%% let's make it do it
\makeatletter
\xpatchcmd{\appendix}
  {\par}
  {\addcontentsline{toc}{section}{\@currentlabelname}\par}
  {}{}
\makeatother

%% Disable longtable counter
%% https://tex.stackexchange.com/a/248395/211326

\usepackage{etoolbox}

\makeatletter
\patchcmd{\LT@caption}
  {\bgroup}
  {\bgroup\global\LTpatch@captiontrue}
  {}{}
\patchcmd{\longtable}
  {\par}
  {\par\global\LTpatch@captionfalse}
  {}{}
\apptocmd{\endlongtable}
  {\ifLTpatch@caption\else\addtocounter{table}{-1}\fi}
  {}{}
\newif\ifLTpatch@caption
\makeatother

\begin{document}

\maketitle


\hypertarget{toc}{}
\tableofcontents
\newpage
\section[Introduction]{Tipos de datos}

\setcounter{secnumdepth}{3}

\setlength\LTleft{0pt}




\subsection{🎯 Motivación}\label{motivaciuxf3n}

Los tipos de datos son la base de la programación. Determinan qué tipo
de información puede almacenar una variable, qué operaciones puedes
realizar con ella y cómo se comporta en diferentes situaciones. Entender
los tipos de datos te permite escribir código más eficiente, evitar
errores inesperados y aprovechar al máximo las capacidades de Python.

\subsection{💡 Intuición}\label{intuiciuxf3n}

Imagina que tienes diferentes tipos de contenedores:

\begin{itemize}
\tightlist
\item
  \textbf{Números} son como una calculadora: puedes sumarlos, restarlos,
  multiplicarlos.
\item
  \textbf{Texto (strings)} son como palabras: puedes unirlas,
  dividirlas, buscar letras.
\item
  \textbf{Listas} son como una bandeja con compartimentos: puedes
  guardar varios elementos ordenados.
\item
  \textbf{Booleanos} son como un interruptor: solo pueden estar
  encendidos (True) o apagados (False).
\end{itemize}

Cada contenedor tiene sus propias reglas y operaciones, y Python
necesita saber con qué tipo está trabajando para actuar correctamente.

\subsection{🧠 Modelo mental}\label{modelo-mental}

En Python:

\begin{enumerate}
\def\labelenumi{\arabic{enumi}.}
\tightlist
\item
  \textbf{Tipado dinámico}: no declaras el tipo; Python lo deduce al
  asignar.
\item
  \textbf{Todo es un objeto}: incluso los números y el texto son objetos
  con métodos.
\item
  \textbf{Los tipos determinan el comportamiento}: lo que puedes hacer
  con una variable depende de su tipo.
\item
  \textbf{\texttt{type()}}: tu herramienta para verificar el tipo de
  cualquier variable.
\end{enumerate}

\subsection{📐 Desarrollo}\label{desarrollo}

\subsubsection{1. Tipos de datos
principales}\label{tipos-de-datos-principales}

{\def\LTcaptype{none} % do not increment counter
\begin{longtable}[]{@{}
  >{\raggedright\arraybackslash}p{(\linewidth - 6\tabcolsep) * \real{0.1250}}
  >{\raggedright\arraybackslash}p{(\linewidth - 6\tabcolsep) * \real{0.3036}}
  >{\raggedright\arraybackslash}p{(\linewidth - 6\tabcolsep) * \real{0.2679}}
  >{\raggedright\arraybackslash}p{(\linewidth - 6\tabcolsep) * \real{0.3036}}@{}}
\toprule\noalign{}
\begin{minipage}[b]{\linewidth}\raggedright
Categoría
\end{minipage} & \begin{minipage}[b]{\linewidth}\raggedright
Tipos
\end{minipage} & \begin{minipage}[b]{\linewidth}\raggedright
Ejemplo
\end{minipage} & \begin{minipage}[b]{\linewidth}\raggedright
Cuándo usar
\end{minipage} \\
\midrule\noalign{}
\endhead
\bottomrule\noalign{}
\endlastfoot
\textbf{Texto} & \texttt{str} & \texttt{"Hola"} & Cualquier cadena de
caracteres \\
\textbf{Números} & \texttt{int}, \texttt{float}, \texttt{complex} &
\texttt{5}, \texttt{3.14}, \texttt{1j} & Operaciones matemáticas \\
\textbf{Booleanos} & \texttt{bool} & \texttt{True}, \texttt{False} &
Condiciones lógicas \\
\textbf{Secuencias} & \texttt{list}, \texttt{tuple}, \texttt{range} &
\texttt{{[}1,2,3{]}}, \texttt{(1,2)}, \texttt{range(5)} & Colecciones de
elementos ordenados \\
\textbf{Mapeo} & \texttt{dict} & \texttt{\{"clave":\ "valor"\}} & Pares
clave-valor \\
\textbf{Conjuntos} & \texttt{set}, \texttt{frozenset} &
\texttt{\{1,2,3\}} & Colecciones sin duplicados \\
\textbf{Binarios} & \texttt{bytes}, \texttt{bytearray},
\texttt{memoryview} & \texttt{b"Hello"} & Datos binarios \\
\textbf{Nulo} & \texttt{NoneType} & \texttt{None} & Valor ``vacío'' o
``sin valor'' \\
\end{longtable}
}

\subsubsection{\texorpdfstring{2. Verificar el tipo con
\texttt{type()}}{2. Verificar el tipo con type()}}\label{verificar-el-tipo-con-type}

\begin{Shaded}
\begin{Highlighting}[]
\NormalTok{x }\OperatorTok{=} \DecValTok{5}
\BuiltInTok{print}\NormalTok{(}\BuiltInTok{type}\NormalTok{(x))  }\CommentTok{\# \textless{}class \textquotesingle{}int\textquotesingle{}\textgreater{}}

\NormalTok{y }\OperatorTok{=} \StringTok{"Hola"}
\BuiltInTok{print}\NormalTok{(}\BuiltInTok{type}\NormalTok{(y))  }\CommentTok{\# \textless{}class \textquotesingle{}str\textquotesingle{}\textgreater{}}

\NormalTok{z }\OperatorTok{=} \FloatTok{3.14}
\BuiltInTok{print}\NormalTok{(}\BuiltInTok{type}\NormalTok{(z))  }\CommentTok{\# \textless{}class \textquotesingle{}float\textquotesingle{}\textgreater{}}

\NormalTok{w }\OperatorTok{=} \VariableTok{True}
\BuiltInTok{print}\NormalTok{(}\BuiltInTok{type}\NormalTok{(w))  }\CommentTok{\# \textless{}class \textquotesingle{}bool\textquotesingle{}\textgreater{}}
\end{Highlighting}
\end{Shaded}

\subsubsection{3. Asignación automática de
tipo}\label{asignaciuxf3n-automuxe1tica-de-tipo}

Python asigna el tipo cuando asignas un valor:

\begin{Shaded}
\begin{Highlighting}[]
\CommentTok{\# Asignación automática}
\NormalTok{nombre }\OperatorTok{=} \StringTok{"Ana"}        \CommentTok{\# str}
\NormalTok{edad }\OperatorTok{=} \DecValTok{25}             \CommentTok{\# int}
\NormalTok{precio }\OperatorTok{=} \FloatTok{99.99}        \CommentTok{\# float}
\NormalTok{es\_valido }\OperatorTok{=} \VariableTok{True}      \CommentTok{\# bool}
\NormalTok{lista }\OperatorTok{=}\NormalTok{ [}\DecValTok{1}\NormalTok{, }\DecValTok{2}\NormalTok{, }\DecValTok{3}\NormalTok{]     }\CommentTok{\# list}
\end{Highlighting}
\end{Shaded}

\subsubsection{4. Casting (conversión
explícita)}\label{casting-conversiuxf3n-expluxedcita}

Puedes convertir entre tipos usando funciones constructoras:

\begin{Shaded}
\begin{Highlighting}[]
\CommentTok{\# Convertir a string}
\BuiltInTok{str}\NormalTok{(}\DecValTok{25}\NormalTok{)        }\CommentTok{\# "25"}
\BuiltInTok{str}\NormalTok{(}\FloatTok{3.14}\NormalTok{)      }\CommentTok{\# "3.14"}

\CommentTok{\# Convertir a entero}
\BuiltInTok{int}\NormalTok{(}\StringTok{"25"}\NormalTok{)      }\CommentTok{\# 25}
\BuiltInTok{int}\NormalTok{(}\FloatTok{3.14}\NormalTok{)      }\CommentTok{\# 3 (trunca decimales)}

\CommentTok{\# Convertir a float}
\BuiltInTok{float}\NormalTok{(}\StringTok{"3.14"}\NormalTok{)  }\CommentTok{\# 3.14}
\BuiltInTok{float}\NormalTok{(}\DecValTok{5}\NormalTok{)       }\CommentTok{\# 5.0}

\CommentTok{\# Convertir a booleano}
\BuiltInTok{bool}\NormalTok{(}\DecValTok{0}\NormalTok{)        }\CommentTok{\# False (cero es False)}
\BuiltInTok{bool}\NormalTok{(}\DecValTok{1}\NormalTok{)        }\CommentTok{\# True}
\BuiltInTok{bool}\NormalTok{(}\StringTok{""}\NormalTok{)       }\CommentTok{\# False (string vacío es False)}
\BuiltInTok{bool}\NormalTok{(}\StringTok{"Hola"}\NormalTok{)   }\CommentTok{\# True}
\end{Highlighting}
\end{Shaded}

\subsubsection{5. Verificar el tipo en
código}\label{verificar-el-tipo-en-cuxf3digo}

A veces necesitas saber el tipo para tomar decisiones:

\begin{Shaded}
\begin{Highlighting}[]
\NormalTok{dato }\OperatorTok{=} \DecValTok{25}

\ControlFlowTok{if} \BuiltInTok{type}\NormalTok{(dato) }\OperatorTok{==} \BuiltInTok{int}\NormalTok{:}
    \BuiltInTok{print}\NormalTok{(}\StringTok{"Es un entero"}\NormalTok{)}
\ControlFlowTok{elif} \BuiltInTok{type}\NormalTok{(dato) }\OperatorTok{==} \BuiltInTok{str}\NormalTok{:}
    \BuiltInTok{print}\NormalTok{(}\StringTok{"Es texto"}\NormalTok{)}

\CommentTok{\# Alternativa más idiomática}
\ControlFlowTok{if} \BuiltInTok{isinstance}\NormalTok{(dato, }\BuiltInTok{int}\NormalTok{):}
    \BuiltInTok{print}\NormalTok{(}\StringTok{"Es un entero"}\NormalTok{)}
\end{Highlighting}
\end{Shaded}

\subsection{🔬 Ejemplos}\label{ejemplos}

\textbf{Mínimo ---} Tipos básicos:

\begin{Shaded}
\begin{Highlighting}[]
\NormalTok{texto }\OperatorTok{=} \StringTok{"Python"}          \CommentTok{\# str}
\NormalTok{entero }\OperatorTok{=} \DecValTok{42}              \CommentTok{\# int}
\NormalTok{decimal }\OperatorTok{=} \FloatTok{3.14}           \CommentTok{\# float}
\NormalTok{booleano }\OperatorTok{=} \VariableTok{True}          \CommentTok{\# bool}
\NormalTok{none }\OperatorTok{=} \VariableTok{None}              \CommentTok{\# NoneType}

\BuiltInTok{print}\NormalTok{(}\BuiltInTok{type}\NormalTok{(texto))    }\CommentTok{\# \textless{}class \textquotesingle{}str\textquotesingle{}\textgreater{}}
\BuiltInTok{print}\NormalTok{(}\BuiltInTok{type}\NormalTok{(entero))   }\CommentTok{\# \textless{}class \textquotesingle{}int\textquotesingle{}\textgreater{}}
\end{Highlighting}
\end{Shaded}

\textbf{Progresivo ---} Conversión entre tipos:

\begin{Shaded}
\begin{Highlighting}[]
\CommentTok{\# Conversión de tipos}
\NormalTok{edad }\OperatorTok{=} \StringTok{"25"}
\BuiltInTok{print}\NormalTok{(}\BuiltInTok{type}\NormalTok{(edad))      }\CommentTok{\# str}
\BuiltInTok{print}\NormalTok{(}\BuiltInTok{type}\NormalTok{(}\BuiltInTok{int}\NormalTok{(edad))) }\CommentTok{\# int}

\CommentTok{\# Problemas comunes}
\CommentTok{\# print("Edad: " + 25)  \# ❌ Error}
\BuiltInTok{print}\NormalTok{(}\StringTok{"Edad: "} \OperatorTok{+} \BuiltInTok{str}\NormalTok{(}\DecValTok{25}\NormalTok{))  }\CommentTok{\# ✅ Convierte a string}

\CommentTok{\# Operaciones con números}
\NormalTok{x }\OperatorTok{=} \DecValTok{5}
\NormalTok{y }\OperatorTok{=} \DecValTok{2}
\BuiltInTok{print}\NormalTok{(x }\OperatorTok{/}\NormalTok{ y)   }\CommentTok{\# 2.5 (float)}
\BuiltInTok{print}\NormalTok{(x }\OperatorTok{//}\NormalTok{ y)  }\CommentTok{\# 2 (int, división entera)}
\BuiltInTok{print}\NormalTok{(x }\OperatorTok{\%}\NormalTok{ y)   }\CommentTok{\# 1 (resto)}
\end{Highlighting}
\end{Shaded}

\textbf{Real ---} Sistema de usuario con validación:

\begin{Shaded}
\begin{Highlighting}[]
\CommentTok{\# Simulación de entrada de usuario}
\NormalTok{nombre }\OperatorTok{=} \StringTok{"Carlos"}        \CommentTok{\# str}
\NormalTok{edad }\OperatorTok{=} \DecValTok{28}               \CommentTok{\# int}
\NormalTok{altura }\OperatorTok{=} \FloatTok{1.75}           \CommentTok{\# float}
\NormalTok{es\_estudiante }\OperatorTok{=} \VariableTok{True}    \CommentTok{\# bool}

\CommentTok{\# Validaciones}
\BuiltInTok{print}\NormalTok{(}\SpecialStringTok{f"Tipo de nombre: }\SpecialCharTok{\{}\BuiltInTok{type}\NormalTok{(nombre)}\SpecialCharTok{\}}\SpecialStringTok{"}\NormalTok{)}
\BuiltInTok{print}\NormalTok{(}\SpecialStringTok{f"Tipo de edad: }\SpecialCharTok{\{}\BuiltInTok{type}\NormalTok{(edad)}\SpecialCharTok{\}}\SpecialStringTok{"}\NormalTok{)}
\BuiltInTok{print}\NormalTok{(}\SpecialStringTok{f"Tipo de altura: }\SpecialCharTok{\{}\BuiltInTok{type}\NormalTok{(altura)}\SpecialCharTok{\}}\SpecialStringTok{"}\NormalTok{)}
\BuiltInTok{print}\NormalTok{(}\SpecialStringTok{f"Tipo de estudiante: }\SpecialCharTok{\{}\BuiltInTok{type}\NormalTok{(es\_estudiante)}\SpecialCharTok{\}}\SpecialStringTok{"}\NormalTok{)}

\CommentTok{\# Conversiones necesarias}
\NormalTok{edad\_texto }\OperatorTok{=} \BuiltInTok{str}\NormalTok{(edad)}
\NormalTok{altura\_texto }\OperatorTok{=} \BuiltInTok{str}\NormalTok{(altura)}

\BuiltInTok{print}\NormalTok{(}\SpecialStringTok{f"Edad: }\SpecialCharTok{\{}\NormalTok{edad\_texto}\SpecialCharTok{\}}\SpecialStringTok{ años"}\NormalTok{)}
\BuiltInTok{print}\NormalTok{(}\SpecialStringTok{f"Altura: }\SpecialCharTok{\{}\NormalTok{altura\_texto}\SpecialCharTok{\}}\SpecialStringTok{ metros"}\NormalTok{)}
\end{Highlighting}
\end{Shaded}

\subsection{⚠️ Errores comunes}\label{errores-comunes}

{\def\LTcaptype{none} % do not increment counter
\begin{longtable}[]{@{}
  >{\raggedright\arraybackslash}p{(\linewidth - 4\tabcolsep) * \real{0.3305}}
  >{\raggedright\arraybackslash}p{(\linewidth - 4\tabcolsep) * \real{0.3136}}
  >{\raggedright\arraybackslash}p{(\linewidth - 4\tabcolsep) * \real{0.3559}}@{}}
\toprule\noalign{}
\begin{minipage}[b]{\linewidth}\raggedright
Error
\end{minipage} & \begin{minipage}[b]{\linewidth}\raggedright
Síntoma
\end{minipage} & \begin{minipage}[b]{\linewidth}\raggedright
Solución
\end{minipage} \\
\midrule\noalign{}
\endhead
\bottomrule\noalign{}
\endlastfoot
\textbf{Sumar string + número} & \texttt{TypeError} & Convierte:
\texttt{str(edad)} o usa coma en print \\
\textbf{Dividir por cero} & \texttt{ZeroDivisionError} & Verifica que el
divisor no sea 0 \\
\textbf{Índice fuera de rango} & \texttt{IndexError} & Verifica la
longitud de la lista \\
\textbf{Comparar tipos incompatibles} & \texttt{TypeError} & Convierte
al mismo tipo antes de comparar \\
\textbf{Usar \texttt{==} en lugar de \texttt{is} con None} & Comparación
válida pero no idiomática & Usa \texttt{is\ None} para verificar None \\
\end{longtable}
}

\subsection{✅ Buenas prácticas}\label{buenas-pruxe1cticas}

\begin{enumerate}
\def\labelenumi{\arabic{enumi}.}
\item
  \textbf{Usa \texttt{isinstance()} en lugar de \texttt{type()} para
  verificar tipos}:

\begin{Shaded}
\begin{Highlighting}[]
\CommentTok{\# ✅ Más flexible}
\ControlFlowTok{if} \BuiltInTok{isinstance}\NormalTok{(valor, }\BuiltInTok{int}\NormalTok{):}
    \BuiltInTok{print}\NormalTok{(}\StringTok{"Es entero"}\NormalTok{)}

\CommentTok{\# ❌ Menos flexible (no maneja herencia)}
\ControlFlowTok{if} \BuiltInTok{type}\NormalTok{(valor) }\OperatorTok{==} \BuiltInTok{int}\NormalTok{:}
    \BuiltInTok{print}\NormalTok{(}\StringTok{"Es entero"}\NormalTok{)}
\end{Highlighting}
\end{Shaded}
\item
  \textbf{Convierte tipos explícitamente cuando sea necesario}:

\begin{Shaded}
\begin{Highlighting}[]
\NormalTok{edad }\OperatorTok{=} \BuiltInTok{input}\NormalTok{(}\StringTok{"Edad: "}\NormalTok{)  }\CommentTok{\# input siempre devuelve str}
\NormalTok{edad\_int }\OperatorTok{=} \BuiltInTok{int}\NormalTok{(edad)    }\CommentTok{\# Convierte a int}
\end{Highlighting}
\end{Shaded}
\item
  \textbf{Sé consciente de los tipos en operaciones mixtas}:

\begin{Shaded}
\begin{Highlighting}[]
\CommentTok{\# La división siempre da float}
\BuiltInTok{print}\NormalTok{(}\DecValTok{5} \OperatorTok{/} \DecValTok{2}\NormalTok{)   }\CommentTok{\# 2.5 (float)}
\BuiltInTok{print}\NormalTok{(}\DecValTok{5} \OperatorTok{//} \DecValTok{2}\NormalTok{)  }\CommentTok{\# 2 (int)}

\CommentTok{\# Suma de int + float da float}
\BuiltInTok{print}\NormalTok{(}\DecValTok{5} \OperatorTok{+} \FloatTok{2.0}\NormalTok{) }\CommentTok{\# 7.0 (float)}
\end{Highlighting}
\end{Shaded}
\item
  \textbf{Usa nombres de variables que indiquen el tipo}
  (implícitamente):

\begin{Shaded}
\begin{Highlighting}[]
\CommentTok{\# Convención}
\NormalTok{nombres\_lista }\OperatorTok{=}\NormalTok{ [}\StringTok{"Ana"}\NormalTok{, }\StringTok{"Carlos"}\NormalTok{]}
\NormalTok{contador\_int }\OperatorTok{=} \DecValTok{5}
\NormalTok{es\_valido\_bool }\OperatorTok{=} \VariableTok{True}
\end{Highlighting}
\end{Shaded}
\item
  \textbf{Documenta los tipos en funciones} (type hints):

\begin{Shaded}
\begin{Highlighting}[]
\KeywordTok{def}\NormalTok{ saludar(nombre: }\BuiltInTok{str}\NormalTok{) }\OperatorTok{{-}\textgreater{}} \BuiltInTok{str}\NormalTok{:}
    \ControlFlowTok{return} \SpecialStringTok{f"Hola, }\SpecialCharTok{\{}\NormalTok{nombre}\SpecialCharTok{\}}\SpecialStringTok{"}
\end{Highlighting}
\end{Shaded}
\end{enumerate}

\subsection{📝 Ejercicio}\label{ejercicio}

\textbf{Básico} --- Exploración de tipos

Escribe un programa que:

\begin{enumerate}
\def\labelenumi{\arabic{enumi}.}
\tightlist
\item
  Cree variables de los siguientes tipos: \texttt{str}, \texttt{int},
  \texttt{float}, \texttt{bool}, \texttt{list}, \texttt{dict},
  \texttt{NoneType}
\item
  Muestre el tipo de cada variable con \texttt{type()}
\item
  Convierta las siguientes variables:

  \begin{itemize}
  \tightlist
  \item
    El \texttt{str} a \texttt{int}
  \item
    El \texttt{int} a \texttt{float}
  \item
    El \texttt{float} a \texttt{str}
  \item
    El \texttt{bool} a \texttt{int}
  \end{itemize}
\item
  Muestre los resultados de cada conversión y sus tipos
\end{enumerate}

\begin{center}\rule{0.5\linewidth}{0.5pt}\end{center}

\textbf{Solución (no mires hasta haberlo intentado):}

\begin{Shaded}
\begin{Highlighting}[]
\CommentTok{\# 1. Variables de diferentes tipos}
\NormalTok{texto }\OperatorTok{=} \StringTok{"42"}
\NormalTok{numero\_entero }\OperatorTok{=} \DecValTok{25}
\NormalTok{numero\_decimal }\OperatorTok{=} \FloatTok{3.14}
\NormalTok{logico }\OperatorTok{=} \VariableTok{True}
\NormalTok{lista\_ejemplo }\OperatorTok{=}\NormalTok{ [}\DecValTok{1}\NormalTok{, }\DecValTok{2}\NormalTok{, }\DecValTok{3}\NormalTok{]}
\NormalTok{diccionario\_ejemplo }\OperatorTok{=}\NormalTok{ \{}\StringTok{"nombre"}\NormalTok{: }\StringTok{"Ana"}\NormalTok{\}}
\NormalTok{nulo }\OperatorTok{=} \VariableTok{None}

\CommentTok{\# 2. Mostrar tipos originales}
\BuiltInTok{print}\NormalTok{(}\StringTok{"Tipos originales:"}\NormalTok{)}
\BuiltInTok{print}\NormalTok{(}\SpecialStringTok{f"texto: }\SpecialCharTok{\{}\BuiltInTok{type}\NormalTok{(texto)}\SpecialCharTok{\}}\SpecialStringTok{"}\NormalTok{)}
\BuiltInTok{print}\NormalTok{(}\SpecialStringTok{f"numero\_entero: }\SpecialCharTok{\{}\BuiltInTok{type}\NormalTok{(numero\_entero)}\SpecialCharTok{\}}\SpecialStringTok{"}\NormalTok{)}
\BuiltInTok{print}\NormalTok{(}\SpecialStringTok{f"numero\_decimal: }\SpecialCharTok{\{}\BuiltInTok{type}\NormalTok{(numero\_decimal)}\SpecialCharTok{\}}\SpecialStringTok{"}\NormalTok{)}
\BuiltInTok{print}\NormalTok{(}\SpecialStringTok{f"logico: }\SpecialCharTok{\{}\BuiltInTok{type}\NormalTok{(logico)}\SpecialCharTok{\}}\SpecialStringTok{"}\NormalTok{)}
\BuiltInTok{print}\NormalTok{(}\SpecialStringTok{f"lista: }\SpecialCharTok{\{}\BuiltInTok{type}\NormalTok{(lista\_ejemplo)}\SpecialCharTok{\}}\SpecialStringTok{"}\NormalTok{)}
\BuiltInTok{print}\NormalTok{(}\SpecialStringTok{f"diccionario: }\SpecialCharTok{\{}\BuiltInTok{type}\NormalTok{(diccionario\_ejemplo)}\SpecialCharTok{\}}\SpecialStringTok{"}\NormalTok{)}
\BuiltInTok{print}\NormalTok{(}\SpecialStringTok{f"nulo: }\SpecialCharTok{\{}\BuiltInTok{type}\NormalTok{(nulo)}\SpecialCharTok{\}}\SpecialStringTok{"}\NormalTok{)}

\BuiltInTok{print}\NormalTok{(}\StringTok{"}\CharTok{\textbackslash{}n}\StringTok{"} \OperatorTok{+} \StringTok{"{-}"} \OperatorTok{*} \DecValTok{40}\NormalTok{)}

\CommentTok{\# 3. Conversiones}
\NormalTok{texto\_a\_int }\OperatorTok{=} \BuiltInTok{int}\NormalTok{(texto)}
\NormalTok{entero\_a\_float }\OperatorTok{=} \BuiltInTok{float}\NormalTok{(numero\_entero)}
\NormalTok{decimal\_a\_str }\OperatorTok{=} \BuiltInTok{str}\NormalTok{(numero\_decimal)}
\NormalTok{bool\_a\_int }\OperatorTok{=} \BuiltInTok{int}\NormalTok{(logico)}

\CommentTok{\# 4. Mostrar resultados}
\BuiltInTok{print}\NormalTok{(}\StringTok{"Resultados de conversiones:"}\NormalTok{)}
\BuiltInTok{print}\NormalTok{(}\SpecialStringTok{f"\textquotesingle{}}\SpecialCharTok{\{}\NormalTok{texto}\SpecialCharTok{\}}\SpecialStringTok{\textquotesingle{} → }\SpecialCharTok{\{}\NormalTok{texto\_a\_int}\SpecialCharTok{\}}\SpecialStringTok{ (tipo: }\SpecialCharTok{\{}\BuiltInTok{type}\NormalTok{(texto\_a\_int)}\SpecialCharTok{\}}\SpecialStringTok{)"}\NormalTok{)}
\BuiltInTok{print}\NormalTok{(}\SpecialStringTok{f"}\SpecialCharTok{\{}\NormalTok{numero\_entero}\SpecialCharTok{\}}\SpecialStringTok{ → }\SpecialCharTok{\{}\NormalTok{entero\_a\_float}\SpecialCharTok{\}}\SpecialStringTok{ (tipo: }\SpecialCharTok{\{}\BuiltInTok{type}\NormalTok{(entero\_a\_float)}\SpecialCharTok{\}}\SpecialStringTok{)"}\NormalTok{)}
\BuiltInTok{print}\NormalTok{(}\SpecialStringTok{f"}\SpecialCharTok{\{}\NormalTok{numero\_decimal}\SpecialCharTok{\}}\SpecialStringTok{ → \textquotesingle{}}\SpecialCharTok{\{}\NormalTok{decimal\_a\_str}\SpecialCharTok{\}}\SpecialStringTok{\textquotesingle{} (tipo: }\SpecialCharTok{\{}\BuiltInTok{type}\NormalTok{(decimal\_a\_str)}\SpecialCharTok{\}}\SpecialStringTok{)"}\NormalTok{)}
\BuiltInTok{print}\NormalTok{(}\SpecialStringTok{f"}\SpecialCharTok{\{}\NormalTok{logico}\SpecialCharTok{\}}\SpecialStringTok{ → }\SpecialCharTok{\{}\NormalTok{bool\_a\_int}\SpecialCharTok{\}}\SpecialStringTok{ (tipo: }\SpecialCharTok{\{}\BuiltInTok{type}\NormalTok{(bool\_a\_int)}\SpecialCharTok{\}}\SpecialStringTok{)"}\NormalTok{)}
\end{Highlighting}
\end{Shaded}

\subsection{🗂️ Síntesis}\label{suxedntesis}

\begin{itemize}
\tightlist
\item
  \textbf{Tipos principales}: \texttt{str}, \texttt{int},
  \texttt{float}, \texttt{bool}, \texttt{list}, \texttt{tuple},
  \texttt{dict}, \texttt{set}, \texttt{NoneType}
\item
  \textbf{\texttt{type()}}: verifica el tipo de una variable
\item
  \textbf{Tipado dinámico}: Python infiere el tipo automáticamente
\item
  \textbf{Casting}: conversión explícita con funciones constructoras
\item
  \textbf{Operaciones}: dependen del tipo de dato
\item
  \textbf{\texttt{isinstance()}}: mejor que \texttt{type()} para
  verificar tipos en herencia
\item
  \textbf{Type hints}: documentación opcional de tipos (Python 3.5+)
\end{itemize}

\subsection{🔗 Qué sigue}\label{quuxe9-sigue}

\begin{itemize}
\tightlist
\item
  \textbf{Relacionado:} variables (asignación y alcance), strings y
  métodos de string
\item
  \textbf{Siguiente tema:} Strings --- aprenderás a trabajar con texto
  en Python, incluyendo métodos, formateo y manipulación avanzada.
\end{itemize}

\section{Publicaciones Similares}\label{publicaciones-similares}

Si te interesó este artículo, te recomendamos que explores otros blogs y
recursos relacionados que pueden ampliar tus conocimientos. Aquí te dejo
algunas sugerencias:

\begin{enumerate}
\def\labelenumi{\arabic{enumi}.}
\tightlist
\item
  \href{https://numerus-scriptum.netlify.app/python/2020-06-19-instalacion-de-anaconda/index.pdf}{\faIcon{file-pdf}}
  \href{https://numerus-scriptum.netlify.app/python/2020-06-19-instalacion-de-anaconda}{Instalacion
  De Anaconda}
\item
  \href{https://numerus-scriptum.netlify.app/python/2020-06-20-configurar-entorno-virtual-python-anaconda/index.pdf}{\faIcon{file-pdf}}
  \href{https://numerus-scriptum.netlify.app/python/2020-06-20-configurar-entorno-virtual-python-anaconda}{Configurar
  Entorno Virtual Python Anaconda}
\item
  \href{https://numerus-scriptum.netlify.app/python/2021-04-17-01-introducion-a-la-programacion-con-python/index.pdf}{\faIcon{file-pdf}}
  \href{https://numerus-scriptum.netlify.app/python/2021-04-17-01-introducion-a-la-programacion-con-python}{01
  Introducion A La Programacion Con Python}
\item
  \href{https://numerus-scriptum.netlify.app/python/2021-05-31-02-variables-expresiones-y-statements-con-python/index.pdf}{\faIcon{file-pdf}}
  \href{https://numerus-scriptum.netlify.app/python/2021-05-31-02-variables-expresiones-y-statements-con-python}{02
  Variables Expresiones Y Statements Con Python}
\item
  \href{https://numerus-scriptum.netlify.app/python/2021-06-07-03-objetos-de-python/index.pdf}{\faIcon{file-pdf}}
  \href{https://numerus-scriptum.netlify.app/python/2021-06-07-03-objetos-de-python}{03
  Objetos De Python}
\item
  \href{https://numerus-scriptum.netlify.app/python/2021-06-14-04-ejecucion-condicional-con-python/index.pdf}{\faIcon{file-pdf}}
  \href{https://numerus-scriptum.netlify.app/python/2021-06-14-04-ejecucion-condicional-con-python}{04
  Ejecucion Condicional Con Python}
\item
  \href{https://numerus-scriptum.netlify.app/python/2021-06-21-05-iteraciones-con-python/index.pdf}{\faIcon{file-pdf}}
  \href{https://numerus-scriptum.netlify.app/python/2021-06-21-05-iteraciones-con-python}{05
  Iteraciones Con Python}
\item
  \href{https://numerus-scriptum.netlify.app/python/2021-08-16-06-funciones-con-python/index.pdf}{\faIcon{file-pdf}}
  \href{https://numerus-scriptum.netlify.app/python/2021-08-16-06-funciones-con-python}{06
  Funciones Con Python}
\item
  \href{https://numerus-scriptum.netlify.app/python/2021-08-23-07-dataframes-con-python/index.pdf}{\faIcon{file-pdf}}
  \href{https://numerus-scriptum.netlify.app/python/2021-08-23-07-dataframes-con-python}{07
  Dataframes Con Python}
\item
  \href{https://numerus-scriptum.netlify.app/python/2021-11-29-08-prediccion-y-metrica-de-performance-con-python/index.pdf}{\faIcon{file-pdf}}
  \href{https://numerus-scriptum.netlify.app/python/2021-11-29-08-prediccion-y-metrica-de-performance-con-python}{08
  Prediccion Y Metrica De Performance Con Python}
\item
  \href{https://numerus-scriptum.netlify.app/python/2021-12-06-09-metodos-de-machine-learning-para-clasificacion-con-python/index.pdf}{\faIcon{file-pdf}}
  \href{https://numerus-scriptum.netlify.app/python/2021-12-06-09-metodos-de-machine-learning-para-clasificacion-con-python}{09
  Metodos De Machine Learning Para Clasificacion Con Python}
\item
  \href{https://numerus-scriptum.netlify.app/python/2021-12-13-10-metodos-de-machine-learning-para-regresion-con-python/index.pdf}{\faIcon{file-pdf}}
  \href{https://numerus-scriptum.netlify.app/python/2021-12-13-10-metodos-de-machine-learning-para-regresion-con-python}{10
  Metodos De Machine Learning Para Regresion Con Python}
\item
  \href{https://numerus-scriptum.netlify.app/python/2025-05-10-visualizacion-de-datos-con-python/index.pdf}{\faIcon{file-pdf}}
  \href{https://numerus-scriptum.netlify.app/python/2025-05-10-visualizacion-de-datos-con-python}{Visualizacion
  De Datos Con Python}
\end{enumerate}

Esperamos que encuentres estas publicaciones igualmente interesantes y
útiles. ¡Disfruta de la lectura!






\end{document}
